\chapter{Analyse et spécification des besoins}\label{chap2}
\pagestyle{fancy}
\lhead{Chapitre 2 : Analyse et spécification des besoins}

% =====================================================================
\section*{Introduction}
\addcontentsline{toc}{section}{Introduction}

Ce chapitre précise ce que la plateforme VigiloChain doit faire, sans dire
comment elle le fait. Nous identifions les acteurs et leurs cas d'utilisation,
puis les besoins fonctionnels et non fonctionnels, et nous fixons le contrôle
d'accès par rôles. Les choix de conception et de réalisation sont traités aux
chapitres suivants.

% =====================================================================
\section{Acteurs et cas d'utilisation}

VigiloChain compte deux acteurs authentifiés : l'\textbf{opérateur}, qui
surveille la chaîne du froid et exploite les résultats de la détection, et
l'\textbf{administrateur}, qui gère en plus le modèle, les paramètres et les
comptes (tableau~\ref{tab:acteurs}). Un \textbf{visiteur} non connecté accède
seulement à la page d'accueil, à l'inscription et à la connexion.

\begin{table}[H]
\centering
\small
\renewcommand{\arraystretch}{1.4}
\setlength{\arrayrulewidth}{0.8pt}
\arrayrulecolor{thead}
\caption{Acteurs du système VigiloChain.}
\label{tab:acteurs}
\begin{tabularx}{\textwidth}{@{}>{\bfseries\raggedright\arraybackslash}p{3.0cm} >{\raggedright\arraybackslash}X@{}}
\rowcolor{thead}
\textcolor{white}{\textbf{Acteur}} & \textcolor{white}{\textbf{Rôle}} \\
\hline
Opérateur &
Suit les mesures du frigo, du thermostat et de la météo rejouées une à une
dans le modèle, consulte les alertes et les recommandations à appliquer, et lit
les résultats du modèle et ses explications (SHAP). \\
\rowcolor{trowa}
Administrateur &
Possède tous les droits de l'opérateur. Il importe un jeu de données,
réentraîne le modèle, règle les paramètres, gère les comptes et consulte le
journal d'audit. \\
\hline
\end{tabularx}
\end{table}

Les capteurs ne sont pas des acteurs : la plateforme travaille sur des traces
réelles enregistrées (TON\_IoT) qu'elle rejoue comme un flux. La
\textbf{détection} n'est pas non plus une action de l'opérateur : le modèle
analyse automatiquement chaque mesure, et l'opérateur observe et décide. Tous
les cas supposent une connexion préalable, non représentée dans le diagramme de
la figure~\ref{fig:casutilisation}.

\begin{figure}[H]
\centering
\begin{tikzpicture}[
  uc/.style={ellipse, draw=thead, fill=white, minimum width=4.3cm,
             minimum height=0.82cm, align=center, font=\scriptsize, inner sep=1pt},
  rel/.style={draw=gray!60, thin},
  gen/.style={draw=thead, thick, -{Triangle[open, length=3mm, width=3mm]}},
  stereo/.style={font=\tiny\itshape, text=thead, fill=white, inner sep=1pt},
  albl/.style={font=\footnotesize\bfseries, text=thead},
  sys/.style={draw=thead!80, dashed, rounded corners=3pt, inner sep=9pt},
  acteur/.pic={
    \draw[thead, thick] (0,0.62) circle (0.17);
    \draw[thead, thick] (0,0.45) -- (0,-0.05);
    \draw[thead, thick] (-0.3,0.28) -- (0.3,0.28);
    \draw[thead, thick] (0,-0.05) -- (-0.22,-0.45);
    \draw[thead, thick] (0,-0.05) -- (0.22,-0.45);
  }
]
% Cas d'utilisation de l'opérateur
\node[uc] (uc1) at (0,0)    {Consulter la vue d'ensemble};
\node[uc] (uc2) at (0,-1.2) {Suivre le flux\\en temps réel};
\node[uc] (uc3) at (0,-2.4) {Comparer détection et attaques,\\lire les recommandations};
\node[uc] (uc4) at (0,-3.6) {Consulter les résultats\\des modèles};
\node[uc] (uc5) at (0,-4.8) {Explorer les explications SHAP};
\node[uc] (uc6) at (0,-6.0) {Gérer son profil\\et son mot de passe};
% Cas d'utilisation de l'administrateur
\node[uc] (uc7)  at (5.6,-1.2) {Importer un jeu\\de données (CSV)};
\node[uc] (uc8)  at (5.6,-2.4) {Réentraîner le modèle\\et gérer les entraînements};
\node[uc] (uc9)  at (5.6,-3.6) {Régler les paramètres};
\node[uc] (uc10) at (5.6,-4.8) {Gérer les comptes};
\node[uc] (uc11) at (5.6,-6.0) {Consulter et purger\\le journal d'audit};
\begin{scope}[on background layer]
  \node[sys, fit=(uc1)(uc6)(uc7)(uc11)] (S) {};
\end{scope}
\node[albl, anchor=south west] at (S.north west) {Système VigiloChain};
\pic at (-3.6,-3.0) {acteur};
\node[albl] (oplbl) at (-3.6,-3.68) {Opérateur};
\coordinate (op) at (-3.3,-2.8);
\pic at (9.3,-3.4) {acteur};
\node[albl] (adlbl) at (9.3,-4.08) {Administrateur};
\coordinate (ad) at (9.0,-3.2);
\foreach \i in {uc1,uc2,uc3,uc4,uc5,uc6} \draw[rel] (op) -- (\i);
\foreach \i in {uc7,uc8,uc9,uc10,uc11} \draw[rel] (ad) -- (\i);
\draw[gen] (adlbl.south) |- ([yshift=-7mm]S.south) -| (oplbl.south);
\node[stereo, below=1pt] at ([yshift=-7mm]S.south)
  {généralisation : l'administrateur hérite des cas de l'opérateur};
\end{tikzpicture}
\caption{Diagramme de cas d'utilisation de VigiloChain.}
\label{fig:casutilisation}
\end{figure}

% =====================================================================
\section{Besoins fonctionnels}

Les besoins fonctionnels décrivent ce que chaque acteur peut faire. Ils sont
regroupés en quatre modules, qui correspondent aux trois groupes du menu du
tableau de bord (Surveillance, Analyse du modèle, Administration) et à la
gestion du compte.

\subsection*{Surveillance}
\begin{itemize}
  \item \textbf{BF1} : L'opérateur peut consulter une vue d'ensemble : dernier
        entraînement, modèle retenu, détection par scénario d'attaque et
        historique des entraînements.
  \item \textbf{BF2} : L'opérateur peut lancer, mettre en pause et arrêter un
        flux temps réel qui rejoue les mesures du frigo, du thermostat et de la
        météo une à une dans le modèle, avec le score, les alertes et la
        réalité côte à côte. Il choisit la période rejouée (entraînement,
        validation ou test) et la vitesse.
\end{itemize}

\subsection*{Analyse du modèle}
\begin{itemize}
  \item \textbf{BF3} : L'opérateur peut comparer, dans le temps, les alertes du
        modèle aux attaques réellement injectées (épisodes détectés, manqués,
        fausses alertes, matrice de confusion) et lire, pour chaque alerte, un
        rapport de recommandation : synthèse des faits mesurés, famille
        d'attaque probable, gravité et actions à mener.
  \item \textbf{BF4} : L'opérateur peut consulter les résultats des modèles
        candidats (rappel, précision, F1, délai de détection) et le meilleur
        modèle retenu.
  \item \textbf{BF5} : L'opérateur peut consulter l'importance des variables
        (SHAP), globale ou pour un scénario d'attaque.
\end{itemize}

\subsection*{Administration}
\begin{itemize}
  \item \textbf{BF6} : L'administrateur peut importer un jeu de données au
        format CSV, réentraîner le modèle retenu sur ce jeu ou sur celui du
        projet, suivre la progression en direct et supprimer un entraînement
        de l'historique.
  \item \textbf{BF7} : L'administrateur peut modifier les paramètres de la
        plateforme : précision minimale, fausses alertes tolérées par heure et
        vitesse par défaut du temps réel.
  \item \textbf{BF8} : L'administrateur peut créer, modifier, activer,
        désactiver et supprimer un compte, et réinitialiser un mot de passe.
        Le compte administrateur lui-même n'est pas modifiable depuis
        l'interface.
  \item \textbf{BF9} : L'administrateur peut consulter et purger le journal
        d'audit.
\end{itemize}

\subsection*{Compte utilisateur}
\begin{itemize}
  \item \textbf{BF10} : Un visiteur peut consulter la page d'accueil (chiffres
        clés publics), s'inscrire (compte opérateur) et se connecter.
  \item \textbf{BF11} : Tout utilisateur connecté peut modifier son profil et
        son mot de passe, changer de thème (clair ou sombre) et se déconnecter.
\end{itemize}

% =====================================================================
\section{Besoins non fonctionnels}
\label{sec:bnf}

Les besoins non fonctionnels fixent la qualité attendue de la plateforme
(tableau~\ref{tab:bnf}).

\begin{table}[H]
\centering
\small
\renewcommand{\arraystretch}{1.4}
\setlength{\arrayrulewidth}{0.8pt}
\arrayrulecolor{thead}
\caption{Besoins non fonctionnels de VigiloChain.}
\label{tab:bnf}
\begin{tabularx}{\textwidth}{@{}>{\bfseries}p{1.4cm} >{\bfseries}p{3.2cm} >{\raggedright\arraybackslash}X@{}}
\rowcolor{thead}
\textcolor{white}{\textbf{Code}} & \textcolor{white}{\textbf{Catégorie}} &
\textcolor{white}{\textbf{Description}} \\
\hline
BNF1 & Performance &
Le temps réel doit suivre la vitesse choisie (de 6 à 180 mesures par seconde
dans l'interface) sans accumuler de retard : le tableau de bord interroge l'API
chaque seconde et regroupe l'affichage une fois par image. \\
\rowcolor{trowa}
BNF2 & Sécurité &
Toute route métier exige un jeton JWT valide, à l'exception des chiffres clés
de la page d'accueil ; les mots de passe sont hachés (bcrypt) ; les droits sont
vérifiés côté serveur selon le rôle, et le tableau de bord masque en plus les
pages réservées. \\
BNF3 & Fiabilité de l'évaluation &
Aucune information du test ne doit fuiter vers l'entraînement : découpage par
blocs d'une heure et zone de garde de 20~min. \\
\rowcolor{trowa}
BNF4 & Explicabilité &
Le comportement du modèle doit pouvoir être justifié par les valeurs SHAP,
affichées dans le tableau de bord. \\
BNF5 & Traçabilité &
Les actions d'administration sont enregistrées avec leur auteur, leur cible
et leur horodatage. \\
\rowcolor{trowa}
BNF6 & Ergonomie &
Interface en français, thème clair ou sombre, menu adapté au rôle et
affichage adaptatif (menu en tiroir sur téléphone et tablette). \\
BNF7 & Maintenabilité &
Code organisé en couches et vérifié par des tests automatisés (pytest). \\
\rowcolor{trowa}
BNF8 & Déploiement &
Installation en une commande (\texttt{docker compose up}). \\
\hline
\end{tabularx}
\end{table}

\subsection{Contrôle d'accès par rôles}
\label{sec:rbac}

Le rôle est porté par le jeton JWT et vérifié \textbf{côté serveur}
(tableau~\ref{tab:rbac}). Le compte administrateur est créé au premier
démarrage à partir des variables d'environnement ; l'inscription publique ne
crée que des comptes opérateur. Un opérateur qui appelle une route
d'administration reçoit le code HTTP~403, et qui ouvre une page réservée voit un
message « Accès refusé ».

\begin{table}[H]
\centering
\small
\renewcommand{\arraystretch}{1.4}
\setlength{\arrayrulewidth}{0.8pt}
\arrayrulecolor{thead}
\caption{Périmètre d'accès par rôle.}
\label{tab:rbac}
\begin{tabularx}{\textwidth}{@{}>{\bfseries}p{3.4cm} >{\raggedright\arraybackslash}X@{}}
\rowcolor{thead}
\textcolor{white}{\textbf{Rôle}} & \textcolor{white}{\textbf{Périmètre}} \\
\hline
Administrateur & Tout ce que fait l'opérateur, plus l'entraînement et
l'import de données, les paramètres, les comptes et le journal d'audit. \\
\rowcolor{trowa}
Opérateur & Surveillance (vue d'ensemble, temps réel), analyse du modèle
(résultats, détection et recommandations, SHAP) et profil. Aucun accès à
l'administration. \\
Visiteur non connecté & Page d'accueil, inscription et connexion
uniquement. \\
\hline
\end{tabularx}
\end{table}

% =====================================================================
\section*{Conclusion}
\addcontentsline{toc}{section}{Conclusion}
Ce chapitre a présenté les deux acteurs authentifiés de VigiloChain, le
visiteur non connecté, ses onze besoins fonctionnels répartis en quatre
modules, ses huit besoins non fonctionnels et le contrôle d'accès par rôles.
Ces besoins servent de référence tout au long du rapport : le chapitre de
réalisation les retrouve un à un dans les pages du tableau de bord et les
routes de l'API (tableau de traçabilité).
